% Generated draft: requires mathematical and source review.
\documentclass[11pt]{article}
\usepackage[margin=25mm]{geometry}
\usepackage{fontspec}
\setmainfont{Latin Modern Roman}
\usepackage{amsmath,amssymb,mathtools,unicode-math}
\setmathfont{Latin Modern Math}
\usepackage{enumitem,needspace,xurl,hyperref}
\setlength{\parindent}{0pt}
\setlength{\parskip}{5pt}
\setlength{\emergencystretch}{3em}
\newcommand{\N}{\mathbb{N}}
\newcommand{\Z}{\mathbb{Z}}
\newcommand{\Q}{\mathbb{Q}}
\newcommand{\R}{\mathbb{R}}
\newcommand{\C}{\mathbb{C}}
\newcommand{\F}{\mathbb{F}}
\newcommand{\A}{\mathbb{A}}
\newcommand{\PP}{\mathbb{P}}
\newcommand{\E}{\mathbb{E}}
\newcommand{\eps}{\varepsilon}
\newcommand{\dd}{\,\mathrm{d}}
\begin{document}
\section*{UnsolvedMath ID: 30004112}
\textbf{Source problem code:} catalogue entry
\subsection{Definitions and mathematical statement}
For an endomorphism \(F:G\to G\), write
\[
G^F=\{g\in G:F(g)=g\}.
\]
If \(G^F\) is finite, the conventional ordinary-character interpretation of \(\operatorname{Irr}(G^F)\) is the set of characters of irreducible finite-dimensional complex representations of \(G^F\). In this interpretation, a character bound typically concerns \(\lvert\chi(g)\rvert\). These conventions explain the notation; the supplied record does not specify the ambient algebraic group, the endomorphism \(F\), the character field, or the precise inequalities intended.

The supplied problem asks the following two questions.
\begin{enumerate}
\item Can the character-bound theorem referred to in the source, currently available for split Levi subgroups, be extended to non-split Levi subgroups?
\item When \(G\) is simple and is not of type \(A\), can useful character bounds be obtained for every pair
\[
(g,\chi)\in G^F\times\operatorname{Irr}(G^F)?
\]
\end{enumerate}
The first question refers to a theorem whose statement, hypotheses, and bound are absent from the supplied record. The meanings of split and non-split Levi subgroups, and of simplicity for \(G\), also require the source context. The second question includes every element and every irreducible character, but does not define what constitutes a useful bound or specify its allowed parameters and dependence on them. Thus the record supplies research targets rather than a fully specified quantitative conjecture; these missing qualifications cannot be recovered from the supplied text alone.
\subsection{Short English statement}
The problem asks whether a character-bound theorem for split Levi subgroups extends to non-split Levi subgroups, and whether useful character bounds hold for every element and every irreducible character when \(G\) is simple and not of type \(A\). The referenced theorem, ambient hypotheses, and required strength of the bounds are unspecified.
\subsection{Sources}
\begin{itemize}
\item Supplied local UnsolvedMath registry record, permanent ID \(30004112\), entitled \emph{Character Bounds beyond Split Levi Subgroups}; exact source problem code \texttt{catalogue entry}. The record belongs to the Oberwolfach Reports Open Problems collection and gives proposed year 2019; no individual proposer is identified.
\item Supplied report citation: \emph{Representations of Finite Groups (2020)}, Oberwolfach Reports, pp.~841--895, DOI \texttt{10.4171/owr/2019/14}. Supplied source URL and literature reference, labelled ``Oberwolfach representation theory report'': \url{https://doi.org/10.4171/owr/2019/14}. The background describes Oberwolfach Reports as documenting mathematical workshops and their open-problem discussions. No specific page, speaker, or theorem number is supplied for these questions.
\item Literature assessment dated exactly \texttt{2026-08-22}: the supplied curation labels the problem open and reports that no general resolution was found for extensions to non-split Levi subgroups and uniform character bounds beyond established simple type-\(A\) and restricted cases. This is supplied dated curation, not an independently verified literature assessment; its time-sensitive status and cited claims require independent verification.
\end{itemize}
\end{document}
